%% ============================================================
%% Lecture 6: Multiple Regression Analysis with Qualitative Information
%% Intermediate Econometrics — SUFE
%% Wooldridge, Introductory Econometrics, 8th edition, Chapter 7
%% ============================================================
\documentclass[aspectratio=169, 12pt]{beamer}
% ==============================================================================
% header.tex — LaTeX Preamble for Intermediate Econometrics
% Shandong University of Finance and Economics (SUFE)
% Textbook: Wooldridge, Introductory Econometrics, 8th edition
% ==============================================================================
%
% USAGE (from Slides/ directory):
%   % ==============================================================================
% header.tex — LaTeX Preamble for Intermediate Econometrics
% Shandong University of Finance and Economics (SUFE)
% Textbook: Wooldridge, Introductory Econometrics, 8th edition
% ==============================================================================
%
% USAGE (from Slides/ directory):
%   % ==============================================================================
% header.tex — LaTeX Preamble for Intermediate Econometrics
% Shandong University of Finance and Economics (SUFE)
% Textbook: Wooldridge, Introductory Econometrics, 8th edition
% ==============================================================================
%
% USAGE (from Slides/ directory):
%   \input{header}          % in Beamer .tex files (TEXINPUTS=../Preambles:)
%
% COMPILE (3-pass XeLaTeX):
%   cd Slides
%   TEXINPUTS=../Preambles:$TEXINPUTS xelatex -interaction=nonstopmode file.tex
%   BIBINPUTS=..:$BIBINPUTS bibtex file
%   TEXINPUTS=../Preambles:$TEXINPUTS xelatex -interaction=nonstopmode file.tex
%   TEXINPUTS=../Preambles:$TEXINPUTS xelatex -interaction=nonstopmode file.tex
% ==============================================================================


% --- Essential packages -------------------------------------------------------
\usepackage{fontspec}           % XeLaTeX font support (must load early)
\usepackage{amsmath}            % Math environments (align, equation, etc.)
\usepackage{amssymb}            % Math symbols (mathbb, etc.)
\usepackage{bm}                 % Bold math (\bm{})
\usepackage{mathtools}          % Extended math tools (coloneqq, etc.)


% --- Table packages -----------------------------------------------------------
\usepackage{booktabs}           % Professional tables (\toprule, \midrule, \bottomrule)
\usepackage{threeparttable}     % Table notes (tablenotes environment)
\usepackage{siunitx}            % Number alignment in tables (S column type)
\sisetup{
  table-format = 1.4,
  table-number-alignment = center,
  detect-weight = true,
  detect-inline-weight = math
}


% --- Figure and graphics packages --------------------------------------------
\usepackage{graphicx}           % Include figures (\includegraphics)
\usepackage{subcaption}         % Subfigures (subfigure environment)
\usepackage{tikz}               % TikZ diagrams
\usepackage{pgfplots}           % Data plots
\pgfplotsset{compat=1.18}
\usetikzlibrary{arrows.meta, positioning, shapes.geometric,
                decorations.pathreplacing, calc, patterns}


% --- Color and boxes ----------------------------------------------------------
\usepackage{xcolor}             % Extended color support
\usepackage{tcolorbox}          % Custom box environments
\tcbuselibrary{skins, breakable, theorems}


% --- Other utility packages ---------------------------------------------------
\usepackage{hyperref}           % Hyperlinks
\usepackage{natbib}             % Bibliography (\citet, \citep, \citealt)
\usepackage{caption}            % Figure/table captions
\usepackage{appendixnumberbeamer} % Appendix frame numbering


% ==============================================================================
% SUFE INSTITUTIONAL COLORS
% ==============================================================================
% Update hex values to match official SUFE brand colors if available.

\definecolor{SUFEBlue}{HTML}{012169}     % Primary: deep navy
\definecolor{SUFEGold}{HTML}{B9975B}     % Secondary: warm gold
\definecolor{SUFEYellow}{HTML}{F2A900}   % Highlight: bright gold/yellow
\definecolor{SUFELight}{HTML}{E8EDF5}    % Light background: pale blue-gray
\definecolor{SUFEDark}{HTML}{1a1a1a}     % Body text: near-black
\definecolor{positive}{HTML}{2E7D32}     % Positive/correct: dark green
\definecolor{negative}{HTML}{C62828}     % Negative/incorrect: dark red


% ==============================================================================
% BEAMER THEME & COLOR SETUP
% ==============================================================================

\usetheme{default}
\useinnertheme{default}
\useoutertheme{default}
\usecolortheme{default}

% Frame title
\setbeamercolor{frametitle}{fg=SUFEBlue, bg=white}
\setbeamerfont{frametitle}{size=\large, series=\bfseries}
\setbeamertemplate{frametitle}[default][left]
\addtobeamertemplate{frametitle}{\vspace{0.05em}}{%
  \vspace{0.05em}{\color{SUFEGold}\rule{\textwidth}{0.5pt}}\vspace{-0.1em}}

% Title page
\setbeamercolor{title}{fg=SUFEBlue}
\setbeamercolor{subtitle}{fg=SUFEGold}
\setbeamercolor{author}{fg=SUFEBlue}
\setbeamercolor{institute}{fg=SUFEDark}
\setbeamercolor{date}{fg=SUFEGold}

% Structure elements
\setbeamercolor{structure}{fg=SUFEBlue}
\setbeamercolor{item}{fg=SUFEBlue}
\setbeamercolor{subitem}{fg=SUFEGold}
\setbeamercolor{subsubitem}{fg=SUFEGold}
\setbeamerfont{itemize/enumerate body}{size=\normalsize}
\setbeamerfont{itemize/enumerate subbody}{size=\small}

% Block environments
\setbeamercolor{block title}{fg=white, bg=SUFEBlue}
\setbeamercolor{block body}{fg=SUFEDark, bg=SUFELight}
\setbeamercolor{block title alerted}{fg=white, bg=red!70!black}
\setbeamercolor{block body alerted}{fg=SUFEDark, bg=red!5}
\setbeamercolor{block title example}{fg=white, bg=SUFEGold!80!black}
\setbeamercolor{block body example}{fg=SUFEDark, bg=SUFEGold!8}

% Remove navigation symbols
\setbeamertemplate{navigation symbols}{}

% Footline: author | short title | slide N/N
\setbeamertemplate{footline}{%
  \leavevmode%
  \hbox{%
    \begin{beamercolorbox}[wd=.30\paperwidth,ht=2.25ex,dp=1ex,left,leftskip=2mm]{author in head/foot}%
      \usebeamerfont{author in head/foot}\insertshortauthor
    \end{beamercolorbox}%
    \begin{beamercolorbox}[wd=.40\paperwidth,ht=2.25ex,dp=1ex,center]{title in head/foot}%
      \usebeamerfont{title in head/foot}\insertshorttitle
    \end{beamercolorbox}%
    \begin{beamercolorbox}[wd=.30\paperwidth,ht=2.25ex,dp=1ex,right,rightskip=2mm]{date in head/foot}%
      \usebeamerfont{date in head/foot}%
      \insertframenumber{} / \inserttotalframenumber
    \end{beamercolorbox}}%
  \vskip0pt%
}
\setbeamercolor{author in head/foot}{fg=white, bg=SUFEBlue}
\setbeamercolor{title in head/foot}{fg=white, bg=SUFEBlue!80!black}
\setbeamercolor{date in head/foot}{fg=white, bg=SUFEBlue}

% Section separator slides (large centered section title)
\AtBeginSection[]{
  \begin{frame}
    \centering
    \vfill
    {\Large\bfseries\color{SUFEBlue}\insertsectionhead}
    \vfill
    {\color{SUFEGold}\rule{0.5\textwidth}{1pt}}
    \vfill
  \end{frame}
}


% ==============================================================================
% FONT SETUP (XeLaTeX)
% ==============================================================================
% Using Windows system fonts (Times New Roman / Arial / Consolas).
% On macOS, swap for: TeX Gyre Termes / TeX Gyre Heros / TeX Gyre Cursor
%   or: Times New Roman / Helvetica / Courier New

\setmainfont{Times New Roman}
\setsansfont{Arial}
\setmonofont{Consolas}


% ==============================================================================
% CUSTOM BOX ENVIRONMENTS (tcolorbox)
% ==============================================================================
% Quarto CSS equivalents: .keybox, .highlightbox, .methodbox,
%                         .assumptionbox, .resultbox

% keybox: key definitions, stated assumptions
\newtcolorbox{keybox}{
  enhanced,
  colback  = SUFEGold!10!white,
  colframe = SUFEGold,
  leftrule = 4pt, rightrule = 0pt, toprule = 0pt, bottomrule = 0pt,
  arc = 0pt, outer arc = 0pt,
  boxsep = 0pt, left = 8pt, right = 6pt, top = 4pt, bottom = 4pt,
}

% highlightbox: important theorems, main results
\newtcolorbox{highlightbox}{
  enhanced,
  colback  = SUFEYellow!8!white,
  colframe = SUFEYellow!80!SUFEGold,
  leftrule = 4pt, rightrule = 0pt, toprule = 0pt, bottomrule = 0pt,
  arc = 0pt, outer arc = 0pt,
  boxsep = 0pt, left = 8pt, right = 6pt, top = 4pt, bottom = 4pt,
}

% definitionbox{Title}: formal definitions — OLS, IV, GMM, etc.
\newtcolorbox{definitionbox}[1]{
  enhanced,
  title    = #1,
  colback  = SUFEBlue!4!white,
  colframe = SUFEBlue,
  colbacktitle = SUFEBlue,
  coltitle = white,
  fonttitle= \bfseries\small,
  arc = 4pt, outer arc = 4pt,
  boxsep = 0pt, left = 8pt, right = 6pt, top = 4pt, bottom = 4pt,
  toptitle = 3pt, bottomtitle = 3pt,
}

% examplebox{Title}: empirical examples, Wooldridge data applications
\newtcolorbox{examplebox}[1]{
  enhanced,
  title    = #1,
  colback  = SUFELight,
  colframe = SUFEGold!70!white,
  colbacktitle = SUFEGold!35!white,
  coltitle = SUFEBlue,
  fonttitle= \bfseries\small,
  arc = 4pt, outer arc = 4pt,
  boxsep = 0pt, left = 8pt, right = 6pt, top = 4pt, bottom = 4pt,
  toptitle = 3pt, bottomtitle = 3pt,
}

% assumptionbox: numbered OLS assumptions (MLR.1–MLR.6)
\newtcolorbox{assumptionbox}{
  enhanced,
  colback  = SUFEGold!8!white,
  colframe = SUFEGold,
  arc = 4pt, outer arc = 4pt,
  boxsep = 0pt, left = 8pt, right = 6pt, top = 4pt, bottom = 4pt,
}


% ==============================================================================
% STANDARD MATH MACROS
% ==============================================================================

% --- Operators ----------------------------------------------------------------
\DeclareMathOperator{\E}{\mathbb{E}}       % Expectation
\DeclareMathOperator{\Var}{Var}             % Variance
\DeclareMathOperator{\Cov}{Cov}             % Covariance
\DeclareMathOperator{\Corr}{Corr}           % Correlation
\DeclareMathOperator{\plim}{plim}           % Probability limit
\DeclareMathOperator{\rank}{rank}           % Matrix rank
\DeclareMathOperator{\tr}{tr}               % Matrix trace
\DeclareMathOperator{\se}{se}               % Standard error

% --- Common shortcuts ---------------------------------------------------------
\newcommand{\bhat}{\hat{\beta}}             % OLS point estimator
\newcommand{\bhatt}{\tilde{\beta}}          % Alternative estimator
\newcommand{\eps}{\varepsilon}              % Error term
\newcommand{\epsh}{\hat{\varepsilon}}       % OLS residual
\newcommand{\uh}{\hat{u}}                   % Residual (alternative)
\newcommand{\yh}{\hat{y}}                   % Fitted value

% Matrices and vectors
\newcommand{\Xb}{\mathbf{X}}               % Design matrix
\newcommand{\yb}{\mathbf{y}}               % Outcome vector
\newcommand{\betab}{\boldsymbol{\beta}}    % Parameter vector
\newcommand{\epsb}{\boldsymbol{\varepsilon}} % Error vector

% Goodness-of-fit
\newcommand{\SSR}{\text{SSR}}              % Sum of squared residuals
\newcommand{\SSE}{\text{SSE}}              % Explained sum of squares
\newcommand{\SST}{\text{SST}}              % Total sum of squares
\newcommand{\Rsq}{R^2}                     % R-squared
\newcommand{\aRsq}{\bar{R}^2}             % Adjusted R-squared

% Asymptotic notation
\newcommand{\avar}{\text{Avar}}            % Asymptotic variance
\newcommand{\pto}{\overset{p}{\to}}        % Convergence in probability
\newcommand{\dto}{\overset{d}{\to}}        % Convergence in distribution
\newcommand{\iid}{\overset{\text{iid}}{\sim}}  % i.i.d. distributed
\newcommand{\asim}{\overset{a}{\sim}}      % Approximately distributed

% Econometric estimators
\newcommand{\OLS}{\text{OLS}}
\newcommand{\IV}{\text{IV}}
\newcommand{\TSLS}{\text{2SLS}}
\newcommand{\GMM}{\text{GMM}}
\newcommand{\FE}{\text{FE}}
\newcommand{\RE}{\text{RE}}
\newcommand{\DID}{\text{DiD}}
\newcommand{\RD}{\text{RD}}

% Probability
\newcommand{\prob}{\mathbb{P}}             % Probability
\newcommand{\indep}{\perp\!\!\!\perp}      % Statistical independence

% Distributions
\newcommand{\Normal}{\mathcal{N}}          % Normal distribution
% Usage: X \sim \Normal(\mu, \sigma^2)

% Absolute value and norm
\newcommand{\abs}[1]{\left\lvert #1 \right\rvert}
\newcommand{\norm}[1]{\left\lVert #1 \right\rVert}


% ==============================================================================
% HYPERREF SETUP
% ==============================================================================

\hypersetup{
  colorlinks = true,
  linkcolor  = SUFEBlue,
  citecolor  = SUFEGold,
  urlcolor   = SUFEBlue
}


% ==============================================================================
% BIBLIOGRAPHY SETUP
% ==============================================================================

\bibliographystyle{aer}   % American Economic Review style
% Usage: \bibliography{../Bibliography_base}  (from Slides/ directory)
% In-text: \citet{Wooldridge2020}, \citep{White1980}
          % in Beamer .tex files (TEXINPUTS=../Preambles:)
%
% COMPILE (3-pass XeLaTeX):
%   cd Slides
%   TEXINPUTS=../Preambles:$TEXINPUTS xelatex -interaction=nonstopmode file.tex
%   BIBINPUTS=..:$BIBINPUTS bibtex file
%   TEXINPUTS=../Preambles:$TEXINPUTS xelatex -interaction=nonstopmode file.tex
%   TEXINPUTS=../Preambles:$TEXINPUTS xelatex -interaction=nonstopmode file.tex
% ==============================================================================


% --- Essential packages -------------------------------------------------------
\usepackage{fontspec}           % XeLaTeX font support (must load early)
\usepackage{amsmath}            % Math environments (align, equation, etc.)
\usepackage{amssymb}            % Math symbols (mathbb, etc.)
\usepackage{bm}                 % Bold math (\bm{})
\usepackage{mathtools}          % Extended math tools (coloneqq, etc.)


% --- Table packages -----------------------------------------------------------
\usepackage{booktabs}           % Professional tables (\toprule, \midrule, \bottomrule)
\usepackage{threeparttable}     % Table notes (tablenotes environment)
\usepackage{siunitx}            % Number alignment in tables (S column type)
\sisetup{
  table-format = 1.4,
  table-number-alignment = center,
  detect-weight = true,
  detect-inline-weight = math
}


% --- Figure and graphics packages --------------------------------------------
\usepackage{graphicx}           % Include figures (\includegraphics)
\usepackage{subcaption}         % Subfigures (subfigure environment)
\usepackage{tikz}               % TikZ diagrams
\usepackage{pgfplots}           % Data plots
\pgfplotsset{compat=1.18}
\usetikzlibrary{arrows.meta, positioning, shapes.geometric,
                decorations.pathreplacing, calc, patterns}


% --- Color and boxes ----------------------------------------------------------
\usepackage{xcolor}             % Extended color support
\usepackage{tcolorbox}          % Custom box environments
\tcbuselibrary{skins, breakable, theorems}


% --- Other utility packages ---------------------------------------------------
\usepackage{hyperref}           % Hyperlinks
\usepackage{natbib}             % Bibliography (\citet, \citep, \citealt)
\usepackage{caption}            % Figure/table captions
\usepackage{appendixnumberbeamer} % Appendix frame numbering


% ==============================================================================
% SUFE INSTITUTIONAL COLORS
% ==============================================================================
% Update hex values to match official SUFE brand colors if available.

\definecolor{SUFEBlue}{HTML}{012169}     % Primary: deep navy
\definecolor{SUFEGold}{HTML}{B9975B}     % Secondary: warm gold
\definecolor{SUFEYellow}{HTML}{F2A900}   % Highlight: bright gold/yellow
\definecolor{SUFELight}{HTML}{E8EDF5}    % Light background: pale blue-gray
\definecolor{SUFEDark}{HTML}{1a1a1a}     % Body text: near-black
\definecolor{positive}{HTML}{2E7D32}     % Positive/correct: dark green
\definecolor{negative}{HTML}{C62828}     % Negative/incorrect: dark red


% ==============================================================================
% BEAMER THEME & COLOR SETUP
% ==============================================================================

\usetheme{default}
\useinnertheme{default}
\useoutertheme{default}
\usecolortheme{default}

% Frame title
\setbeamercolor{frametitle}{fg=SUFEBlue, bg=white}
\setbeamerfont{frametitle}{size=\large, series=\bfseries}
\setbeamertemplate{frametitle}[default][left]
\addtobeamertemplate{frametitle}{\vspace{0.05em}}{%
  \vspace{0.05em}{\color{SUFEGold}\rule{\textwidth}{0.5pt}}\vspace{-0.1em}}

% Title page
\setbeamercolor{title}{fg=SUFEBlue}
\setbeamercolor{subtitle}{fg=SUFEGold}
\setbeamercolor{author}{fg=SUFEBlue}
\setbeamercolor{institute}{fg=SUFEDark}
\setbeamercolor{date}{fg=SUFEGold}

% Structure elements
\setbeamercolor{structure}{fg=SUFEBlue}
\setbeamercolor{item}{fg=SUFEBlue}
\setbeamercolor{subitem}{fg=SUFEGold}
\setbeamercolor{subsubitem}{fg=SUFEGold}
\setbeamerfont{itemize/enumerate body}{size=\normalsize}
\setbeamerfont{itemize/enumerate subbody}{size=\small}

% Block environments
\setbeamercolor{block title}{fg=white, bg=SUFEBlue}
\setbeamercolor{block body}{fg=SUFEDark, bg=SUFELight}
\setbeamercolor{block title alerted}{fg=white, bg=red!70!black}
\setbeamercolor{block body alerted}{fg=SUFEDark, bg=red!5}
\setbeamercolor{block title example}{fg=white, bg=SUFEGold!80!black}
\setbeamercolor{block body example}{fg=SUFEDark, bg=SUFEGold!8}

% Remove navigation symbols
\setbeamertemplate{navigation symbols}{}

% Footline: author | short title | slide N/N
\setbeamertemplate{footline}{%
  \leavevmode%
  \hbox{%
    \begin{beamercolorbox}[wd=.30\paperwidth,ht=2.25ex,dp=1ex,left,leftskip=2mm]{author in head/foot}%
      \usebeamerfont{author in head/foot}\insertshortauthor
    \end{beamercolorbox}%
    \begin{beamercolorbox}[wd=.40\paperwidth,ht=2.25ex,dp=1ex,center]{title in head/foot}%
      \usebeamerfont{title in head/foot}\insertshorttitle
    \end{beamercolorbox}%
    \begin{beamercolorbox}[wd=.30\paperwidth,ht=2.25ex,dp=1ex,right,rightskip=2mm]{date in head/foot}%
      \usebeamerfont{date in head/foot}%
      \insertframenumber{} / \inserttotalframenumber
    \end{beamercolorbox}}%
  \vskip0pt%
}
\setbeamercolor{author in head/foot}{fg=white, bg=SUFEBlue}
\setbeamercolor{title in head/foot}{fg=white, bg=SUFEBlue!80!black}
\setbeamercolor{date in head/foot}{fg=white, bg=SUFEBlue}

% Section separator slides (large centered section title)
\AtBeginSection[]{
  \begin{frame}
    \centering
    \vfill
    {\Large\bfseries\color{SUFEBlue}\insertsectionhead}
    \vfill
    {\color{SUFEGold}\rule{0.5\textwidth}{1pt}}
    \vfill
  \end{frame}
}


% ==============================================================================
% FONT SETUP (XeLaTeX)
% ==============================================================================
% Using Windows system fonts (Times New Roman / Arial / Consolas).
% On macOS, swap for: TeX Gyre Termes / TeX Gyre Heros / TeX Gyre Cursor
%   or: Times New Roman / Helvetica / Courier New

\setmainfont{Times New Roman}
\setsansfont{Arial}
\setmonofont{Consolas}


% ==============================================================================
% CUSTOM BOX ENVIRONMENTS (tcolorbox)
% ==============================================================================
% Quarto CSS equivalents: .keybox, .highlightbox, .methodbox,
%                         .assumptionbox, .resultbox

% keybox: key definitions, stated assumptions
\newtcolorbox{keybox}{
  enhanced,
  colback  = SUFEGold!10!white,
  colframe = SUFEGold,
  leftrule = 4pt, rightrule = 0pt, toprule = 0pt, bottomrule = 0pt,
  arc = 0pt, outer arc = 0pt,
  boxsep = 0pt, left = 8pt, right = 6pt, top = 4pt, bottom = 4pt,
}

% highlightbox: important theorems, main results
\newtcolorbox{highlightbox}{
  enhanced,
  colback  = SUFEYellow!8!white,
  colframe = SUFEYellow!80!SUFEGold,
  leftrule = 4pt, rightrule = 0pt, toprule = 0pt, bottomrule = 0pt,
  arc = 0pt, outer arc = 0pt,
  boxsep = 0pt, left = 8pt, right = 6pt, top = 4pt, bottom = 4pt,
}

% definitionbox{Title}: formal definitions — OLS, IV, GMM, etc.
\newtcolorbox{definitionbox}[1]{
  enhanced,
  title    = #1,
  colback  = SUFEBlue!4!white,
  colframe = SUFEBlue,
  colbacktitle = SUFEBlue,
  coltitle = white,
  fonttitle= \bfseries\small,
  arc = 4pt, outer arc = 4pt,
  boxsep = 0pt, left = 8pt, right = 6pt, top = 4pt, bottom = 4pt,
  toptitle = 3pt, bottomtitle = 3pt,
}

% examplebox{Title}: empirical examples, Wooldridge data applications
\newtcolorbox{examplebox}[1]{
  enhanced,
  title    = #1,
  colback  = SUFELight,
  colframe = SUFEGold!70!white,
  colbacktitle = SUFEGold!35!white,
  coltitle = SUFEBlue,
  fonttitle= \bfseries\small,
  arc = 4pt, outer arc = 4pt,
  boxsep = 0pt, left = 8pt, right = 6pt, top = 4pt, bottom = 4pt,
  toptitle = 3pt, bottomtitle = 3pt,
}

% assumptionbox: numbered OLS assumptions (MLR.1–MLR.6)
\newtcolorbox{assumptionbox}{
  enhanced,
  colback  = SUFEGold!8!white,
  colframe = SUFEGold,
  arc = 4pt, outer arc = 4pt,
  boxsep = 0pt, left = 8pt, right = 6pt, top = 4pt, bottom = 4pt,
}


% ==============================================================================
% STANDARD MATH MACROS
% ==============================================================================

% --- Operators ----------------------------------------------------------------
\DeclareMathOperator{\E}{\mathbb{E}}       % Expectation
\DeclareMathOperator{\Var}{Var}             % Variance
\DeclareMathOperator{\Cov}{Cov}             % Covariance
\DeclareMathOperator{\Corr}{Corr}           % Correlation
\DeclareMathOperator{\plim}{plim}           % Probability limit
\DeclareMathOperator{\rank}{rank}           % Matrix rank
\DeclareMathOperator{\tr}{tr}               % Matrix trace
\DeclareMathOperator{\se}{se}               % Standard error

% --- Common shortcuts ---------------------------------------------------------
\newcommand{\bhat}{\hat{\beta}}             % OLS point estimator
\newcommand{\bhatt}{\tilde{\beta}}          % Alternative estimator
\newcommand{\eps}{\varepsilon}              % Error term
\newcommand{\epsh}{\hat{\varepsilon}}       % OLS residual
\newcommand{\uh}{\hat{u}}                   % Residual (alternative)
\newcommand{\yh}{\hat{y}}                   % Fitted value

% Matrices and vectors
\newcommand{\Xb}{\mathbf{X}}               % Design matrix
\newcommand{\yb}{\mathbf{y}}               % Outcome vector
\newcommand{\betab}{\boldsymbol{\beta}}    % Parameter vector
\newcommand{\epsb}{\boldsymbol{\varepsilon}} % Error vector

% Goodness-of-fit
\newcommand{\SSR}{\text{SSR}}              % Sum of squared residuals
\newcommand{\SSE}{\text{SSE}}              % Explained sum of squares
\newcommand{\SST}{\text{SST}}              % Total sum of squares
\newcommand{\Rsq}{R^2}                     % R-squared
\newcommand{\aRsq}{\bar{R}^2}             % Adjusted R-squared

% Asymptotic notation
\newcommand{\avar}{\text{Avar}}            % Asymptotic variance
\newcommand{\pto}{\overset{p}{\to}}        % Convergence in probability
\newcommand{\dto}{\overset{d}{\to}}        % Convergence in distribution
\newcommand{\iid}{\overset{\text{iid}}{\sim}}  % i.i.d. distributed
\newcommand{\asim}{\overset{a}{\sim}}      % Approximately distributed

% Econometric estimators
\newcommand{\OLS}{\text{OLS}}
\newcommand{\IV}{\text{IV}}
\newcommand{\TSLS}{\text{2SLS}}
\newcommand{\GMM}{\text{GMM}}
\newcommand{\FE}{\text{FE}}
\newcommand{\RE}{\text{RE}}
\newcommand{\DID}{\text{DiD}}
\newcommand{\RD}{\text{RD}}

% Probability
\newcommand{\prob}{\mathbb{P}}             % Probability
\newcommand{\indep}{\perp\!\!\!\perp}      % Statistical independence

% Distributions
\newcommand{\Normal}{\mathcal{N}}          % Normal distribution
% Usage: X \sim \Normal(\mu, \sigma^2)

% Absolute value and norm
\newcommand{\abs}[1]{\left\lvert #1 \right\rvert}
\newcommand{\norm}[1]{\left\lVert #1 \right\rVert}


% ==============================================================================
% HYPERREF SETUP
% ==============================================================================

\hypersetup{
  colorlinks = true,
  linkcolor  = SUFEBlue,
  citecolor  = SUFEGold,
  urlcolor   = SUFEBlue
}


% ==============================================================================
% BIBLIOGRAPHY SETUP
% ==============================================================================

\bibliographystyle{aer}   % American Economic Review style
% Usage: \bibliography{../Bibliography_base}  (from Slides/ directory)
% In-text: \citet{Wooldridge2020}, \citep{White1980}
          % in Beamer .tex files (TEXINPUTS=../Preambles:)
%
% COMPILE (3-pass XeLaTeX):
%   cd Slides
%   TEXINPUTS=../Preambles:$TEXINPUTS xelatex -interaction=nonstopmode file.tex
%   BIBINPUTS=..:$BIBINPUTS bibtex file
%   TEXINPUTS=../Preambles:$TEXINPUTS xelatex -interaction=nonstopmode file.tex
%   TEXINPUTS=../Preambles:$TEXINPUTS xelatex -interaction=nonstopmode file.tex
% ==============================================================================


% --- Essential packages -------------------------------------------------------
\usepackage{fontspec}           % XeLaTeX font support (must load early)
\usepackage{amsmath}            % Math environments (align, equation, etc.)
\usepackage{amssymb}            % Math symbols (mathbb, etc.)
\usepackage{bm}                 % Bold math (\bm{})
\usepackage{mathtools}          % Extended math tools (coloneqq, etc.)


% --- Table packages -----------------------------------------------------------
\usepackage{booktabs}           % Professional tables (\toprule, \midrule, \bottomrule)
\usepackage{threeparttable}     % Table notes (tablenotes environment)
\usepackage{siunitx}            % Number alignment in tables (S column type)
\sisetup{
  table-format = 1.4,
  table-number-alignment = center,
  detect-weight = true,
  detect-inline-weight = math
}


% --- Figure and graphics packages --------------------------------------------
\usepackage{graphicx}           % Include figures (\includegraphics)
\usepackage{subcaption}         % Subfigures (subfigure environment)
\usepackage{tikz}               % TikZ diagrams
\usepackage{pgfplots}           % Data plots
\pgfplotsset{compat=1.18}
\usetikzlibrary{arrows.meta, positioning, shapes.geometric,
                decorations.pathreplacing, calc, patterns}


% --- Color and boxes ----------------------------------------------------------
\usepackage{xcolor}             % Extended color support
\usepackage{tcolorbox}          % Custom box environments
\tcbuselibrary{skins, breakable, theorems}


% --- Other utility packages ---------------------------------------------------
\usepackage{hyperref}           % Hyperlinks
\usepackage{natbib}             % Bibliography (\citet, \citep, \citealt)
\usepackage{caption}            % Figure/table captions
\usepackage{appendixnumberbeamer} % Appendix frame numbering


% ==============================================================================
% SUFE INSTITUTIONAL COLORS
% ==============================================================================
% Update hex values to match official SUFE brand colors if available.

\definecolor{SUFEBlue}{HTML}{012169}     % Primary: deep navy
\definecolor{SUFEGold}{HTML}{B9975B}     % Secondary: warm gold
\definecolor{SUFEYellow}{HTML}{F2A900}   % Highlight: bright gold/yellow
\definecolor{SUFELight}{HTML}{E8EDF5}    % Light background: pale blue-gray
\definecolor{SUFEDark}{HTML}{1a1a1a}     % Body text: near-black
\definecolor{positive}{HTML}{2E7D32}     % Positive/correct: dark green
\definecolor{negative}{HTML}{C62828}     % Negative/incorrect: dark red


% ==============================================================================
% BEAMER THEME & COLOR SETUP
% ==============================================================================

\usetheme{default}
\useinnertheme{default}
\useoutertheme{default}
\usecolortheme{default}

% Frame title
\setbeamercolor{frametitle}{fg=SUFEBlue, bg=white}
\setbeamerfont{frametitle}{size=\large, series=\bfseries}
\setbeamertemplate{frametitle}[default][left]
\addtobeamertemplate{frametitle}{\vspace{0.05em}}{%
  \vspace{0.05em}{\color{SUFEGold}\rule{\textwidth}{0.5pt}}\vspace{-0.1em}}

% Title page
\setbeamercolor{title}{fg=SUFEBlue}
\setbeamercolor{subtitle}{fg=SUFEGold}
\setbeamercolor{author}{fg=SUFEBlue}
\setbeamercolor{institute}{fg=SUFEDark}
\setbeamercolor{date}{fg=SUFEGold}

% Structure elements
\setbeamercolor{structure}{fg=SUFEBlue}
\setbeamercolor{item}{fg=SUFEBlue}
\setbeamercolor{subitem}{fg=SUFEGold}
\setbeamercolor{subsubitem}{fg=SUFEGold}
\setbeamerfont{itemize/enumerate body}{size=\normalsize}
\setbeamerfont{itemize/enumerate subbody}{size=\small}

% Block environments
\setbeamercolor{block title}{fg=white, bg=SUFEBlue}
\setbeamercolor{block body}{fg=SUFEDark, bg=SUFELight}
\setbeamercolor{block title alerted}{fg=white, bg=red!70!black}
\setbeamercolor{block body alerted}{fg=SUFEDark, bg=red!5}
\setbeamercolor{block title example}{fg=white, bg=SUFEGold!80!black}
\setbeamercolor{block body example}{fg=SUFEDark, bg=SUFEGold!8}

% Remove navigation symbols
\setbeamertemplate{navigation symbols}{}

% Footline: author | short title | slide N/N
\setbeamertemplate{footline}{%
  \leavevmode%
  \hbox{%
    \begin{beamercolorbox}[wd=.30\paperwidth,ht=2.25ex,dp=1ex,left,leftskip=2mm]{author in head/foot}%
      \usebeamerfont{author in head/foot}\insertshortauthor
    \end{beamercolorbox}%
    \begin{beamercolorbox}[wd=.40\paperwidth,ht=2.25ex,dp=1ex,center]{title in head/foot}%
      \usebeamerfont{title in head/foot}\insertshorttitle
    \end{beamercolorbox}%
    \begin{beamercolorbox}[wd=.30\paperwidth,ht=2.25ex,dp=1ex,right,rightskip=2mm]{date in head/foot}%
      \usebeamerfont{date in head/foot}%
      \insertframenumber{} / \inserttotalframenumber
    \end{beamercolorbox}}%
  \vskip0pt%
}
\setbeamercolor{author in head/foot}{fg=white, bg=SUFEBlue}
\setbeamercolor{title in head/foot}{fg=white, bg=SUFEBlue!80!black}
\setbeamercolor{date in head/foot}{fg=white, bg=SUFEBlue}

% Section separator slides (large centered section title)
\AtBeginSection[]{
  \begin{frame}
    \centering
    \vfill
    {\Large\bfseries\color{SUFEBlue}\insertsectionhead}
    \vfill
    {\color{SUFEGold}\rule{0.5\textwidth}{1pt}}
    \vfill
  \end{frame}
}


% ==============================================================================
% FONT SETUP (XeLaTeX)
% ==============================================================================
% Using Windows system fonts (Times New Roman / Arial / Consolas).
% On macOS, swap for: TeX Gyre Termes / TeX Gyre Heros / TeX Gyre Cursor
%   or: Times New Roman / Helvetica / Courier New

\setmainfont{Times New Roman}
\setsansfont{Arial}
\setmonofont{Consolas}


% ==============================================================================
% CUSTOM BOX ENVIRONMENTS (tcolorbox)
% ==============================================================================
% Quarto CSS equivalents: .keybox, .highlightbox, .methodbox,
%                         .assumptionbox, .resultbox

% keybox: key definitions, stated assumptions
\newtcolorbox{keybox}{
  enhanced,
  colback  = SUFEGold!10!white,
  colframe = SUFEGold,
  leftrule = 4pt, rightrule = 0pt, toprule = 0pt, bottomrule = 0pt,
  arc = 0pt, outer arc = 0pt,
  boxsep = 0pt, left = 8pt, right = 6pt, top = 4pt, bottom = 4pt,
}

% highlightbox: important theorems, main results
\newtcolorbox{highlightbox}{
  enhanced,
  colback  = SUFEYellow!8!white,
  colframe = SUFEYellow!80!SUFEGold,
  leftrule = 4pt, rightrule = 0pt, toprule = 0pt, bottomrule = 0pt,
  arc = 0pt, outer arc = 0pt,
  boxsep = 0pt, left = 8pt, right = 6pt, top = 4pt, bottom = 4pt,
}

% definitionbox{Title}: formal definitions — OLS, IV, GMM, etc.
\newtcolorbox{definitionbox}[1]{
  enhanced,
  title    = #1,
  colback  = SUFEBlue!4!white,
  colframe = SUFEBlue,
  colbacktitle = SUFEBlue,
  coltitle = white,
  fonttitle= \bfseries\small,
  arc = 4pt, outer arc = 4pt,
  boxsep = 0pt, left = 8pt, right = 6pt, top = 4pt, bottom = 4pt,
  toptitle = 3pt, bottomtitle = 3pt,
}

% examplebox{Title}: empirical examples, Wooldridge data applications
\newtcolorbox{examplebox}[1]{
  enhanced,
  title    = #1,
  colback  = SUFELight,
  colframe = SUFEGold!70!white,
  colbacktitle = SUFEGold!35!white,
  coltitle = SUFEBlue,
  fonttitle= \bfseries\small,
  arc = 4pt, outer arc = 4pt,
  boxsep = 0pt, left = 8pt, right = 6pt, top = 4pt, bottom = 4pt,
  toptitle = 3pt, bottomtitle = 3pt,
}

% assumptionbox: numbered OLS assumptions (MLR.1–MLR.6)
\newtcolorbox{assumptionbox}{
  enhanced,
  colback  = SUFEGold!8!white,
  colframe = SUFEGold,
  arc = 4pt, outer arc = 4pt,
  boxsep = 0pt, left = 8pt, right = 6pt, top = 4pt, bottom = 4pt,
}


% ==============================================================================
% STANDARD MATH MACROS
% ==============================================================================

% --- Operators ----------------------------------------------------------------
\DeclareMathOperator{\E}{\mathbb{E}}       % Expectation
\DeclareMathOperator{\Var}{Var}             % Variance
\DeclareMathOperator{\Cov}{Cov}             % Covariance
\DeclareMathOperator{\Corr}{Corr}           % Correlation
\DeclareMathOperator{\plim}{plim}           % Probability limit
\DeclareMathOperator{\rank}{rank}           % Matrix rank
\DeclareMathOperator{\tr}{tr}               % Matrix trace
\DeclareMathOperator{\se}{se}               % Standard error

% --- Common shortcuts ---------------------------------------------------------
\newcommand{\bhat}{\hat{\beta}}             % OLS point estimator
\newcommand{\bhatt}{\tilde{\beta}}          % Alternative estimator
\newcommand{\eps}{\varepsilon}              % Error term
\newcommand{\epsh}{\hat{\varepsilon}}       % OLS residual
\newcommand{\uh}{\hat{u}}                   % Residual (alternative)
\newcommand{\yh}{\hat{y}}                   % Fitted value

% Matrices and vectors
\newcommand{\Xb}{\mathbf{X}}               % Design matrix
\newcommand{\yb}{\mathbf{y}}               % Outcome vector
\newcommand{\betab}{\boldsymbol{\beta}}    % Parameter vector
\newcommand{\epsb}{\boldsymbol{\varepsilon}} % Error vector

% Goodness-of-fit
\newcommand{\SSR}{\text{SSR}}              % Sum of squared residuals
\newcommand{\SSE}{\text{SSE}}              % Explained sum of squares
\newcommand{\SST}{\text{SST}}              % Total sum of squares
\newcommand{\Rsq}{R^2}                     % R-squared
\newcommand{\aRsq}{\bar{R}^2}             % Adjusted R-squared

% Asymptotic notation
\newcommand{\avar}{\text{Avar}}            % Asymptotic variance
\newcommand{\pto}{\overset{p}{\to}}        % Convergence in probability
\newcommand{\dto}{\overset{d}{\to}}        % Convergence in distribution
\newcommand{\iid}{\overset{\text{iid}}{\sim}}  % i.i.d. distributed
\newcommand{\asim}{\overset{a}{\sim}}      % Approximately distributed

% Econometric estimators
\newcommand{\OLS}{\text{OLS}}
\newcommand{\IV}{\text{IV}}
\newcommand{\TSLS}{\text{2SLS}}
\newcommand{\GMM}{\text{GMM}}
\newcommand{\FE}{\text{FE}}
\newcommand{\RE}{\text{RE}}
\newcommand{\DID}{\text{DiD}}
\newcommand{\RD}{\text{RD}}

% Probability
\newcommand{\prob}{\mathbb{P}}             % Probability
\newcommand{\indep}{\perp\!\!\!\perp}      % Statistical independence

% Distributions
\newcommand{\Normal}{\mathcal{N}}          % Normal distribution
% Usage: X \sim \Normal(\mu, \sigma^2)

% Absolute value and norm
\newcommand{\abs}[1]{\left\lvert #1 \right\rvert}
\newcommand{\norm}[1]{\left\lVert #1 \right\rVert}


% ==============================================================================
% HYPERREF SETUP
% ==============================================================================

\hypersetup{
  colorlinks = true,
  linkcolor  = SUFEBlue,
  citecolor  = SUFEGold,
  urlcolor   = SUFEBlue
}


% ==============================================================================
% BIBLIOGRAPHY SETUP
% ==============================================================================

\bibliographystyle{aer}   % American Economic Review style
% Usage: \bibliography{../Bibliography_base}  (from Slides/ directory)
% In-text: \citet{Wooldridge2020}, \citep{White1980}


\title[L6: Qualitative Data]{Multiple Regression Analysis with Qualitative Information}
\subtitle{Intermediate Econometrics}
\author[SUFE Econometrics]{Chengye Jia}
\institute{%
  Department of Economics \\
  Shandong University of Finance and Economics}
\date{Spring 2026 \\ \smallskip \scriptsize Wooldridge, Chapter 7}

%% ============================================================
\begin{document}
%% ============================================================

\begin{frame}[plain]
  \titlepage
\end{frame}

\begin{frame}[shrink=10]{Outline}
  \tableofcontents
\end{frame}

%% ============================================================
\section{Describing Qualitative Information (\S7-1)}
%% ============================================================

\begin{frame}[shrink=20]{Dummy Variables: Definition and Naming}

  Many variables in economics are \textbf{qualitative}, not quantitative.

  \medskip

  \begin{examplebox}{Examples of Qualitative Information}
    \begin{itemize}
      \item Gender: male or female
      \item Marital status: married or single
      \item Employment: employed or unemployed
      \item Education: high school, college, graduate
      \item Computer ownership: yes or no
    \end{itemize}
  \end{examplebox}

  \medskip

  \textbf{Solution:} Encode qualitative info as \textbf{dummy variables} (or \textbf{binary}, \textbf{indicator}, \textbf{zero-one} variables).

  \begin{highlightbox}
    A dummy variable takes value 1 when a condition is \textit{true} and 0 when \textit{false}.

    \textbf{Naming convention:} Use descriptive names that clarify which case = 1.
    \begin{itemize}
      \item Good: \texttt{female} (1 = female, 0 = male)
      \item Bad: \texttt{gender} (unclear which is 1)
      \item Good: \texttt{married} (1 = married, 0 = not married)
      \item Good: \texttt{urban} (1 = urban, 0 = rural)
    \end{itemize}
  \end{highlightbox}

  \medskip

  \begin{keybox}
    \textbf{Why dummies matter:} They allow us to quantify the effect of categorical differences. Instead of analyzing men and women separately, we estimate one model that captures both groups with an intercept shift.
  \end{keybox}

\end{frame}

\begin{frame}[shrink=25]{Data Example: WAGE1 Dataset}

  \textbf{Wooldridge's WAGE1 dataset:} 526 workers from 1976 Current Population Survey

  \medskip

  \begin{center}
    \small
    \begin{tabular}{lllllll}
      \toprule
      \textit{wage} & \textit{educ} & \textit{exper} & \textit{tenure} & \textit{female} & \textit{married} & \textit{nonwhite} \\
      \midrule
      3.10 & 11 & 2 & 0 & 1 & 0 & 0 \\
      3.24 & 12 & 22 & 2 & 1 & 1 & 0 \\
      3.00 & 11 & 2 & 0 & 0 & 0 & 0 \\
      6.00 & 8 & 44 & 28 & 0 & 1 & 0 \\
      5.30 & 12 & 7 & 3 & 0 & 1 & 1 \\
      \vdots & \vdots & \vdots & \vdots & \vdots & \vdots & \vdots \\
      \bottomrule
    \end{tabular}
  \end{center}

  \medskip

  \begin{keybox}
    Note: \texttt{female}, \texttt{married}, \texttt{nonwhite} are dummy variables (0 or 1).
    This dataset allows us to test for gender discrimination, marriage premium, and racial gaps.
  \end{keybox}

\end{frame}

%% ============================================================
\section{A Single Dummy Independent Variable (\S7-2)}
%% ============================================================

\begin{frame}[shrink=25]{A Single Dummy: The Gender-Wage Model}

  \textbf{Question:} Do women earn less than men, \textit{ceteris paribus}?

  \medskip

  \begin{definitionbox}{Simple Model with Dummy (eq.\ 7.1)}
    \[
      \textit{wage} = \beta_0 + \delta_0 \textit{female} + \beta_1 \textit{educ} + u
    \]
    where \texttt{female} = 1 if worker is female, 0 if male.
  \end{definitionbox}

  \medskip

  \textbf{Interpretation:}

  \begin{itemize}
    \item $\beta_0$ = hourly wage for a \textbf{male} with 0 education (intercept; probably not realistic but mathematic necessity)
    \item $\delta_0$ = \textbf{difference in wage} between females and males, holding education fixed (the \textbf{ceteris paribus gender wage gap})
    \item $\beta_1$ = return to one additional year of education (for both men and women; assumes same return for both genders)
  \end{itemize}

  \medskip

  \textbf{Key economic insight:} If $\delta_0 \neq 0$, we have evidence of a systematic wage difference unexplained by education alone. This could reflect discrimination, unobserved skill differences, occupational sorting, or other factors.

  \medskip

  \begin{highlightbox}
    \[
      \E[\textit{wage} \mid \textit{female}, \textit{educ}] =
      \begin{cases}
        \beta_0 + \beta_1 \textit{educ} & \text{if male} \\
        (\beta_0 + \delta_0) + \beta_1 \textit{educ} & \text{if female}
      \end{cases}
    \]

    Geometrically: two \textbf{parallel lines} with the same slope but different intercepts.
  \end{highlightbox}

\end{frame}

\begin{frame}[shrink=15]{Parallel Regression Lines: The Intercept Shift}

  \begin{center}
    \begin{tikzpicture}[scale=1.0]
      % Draw axes
      \draw[->, line width=1pt] (0, 0) -- (14, 0) node[below right] {Education (years)};
      \draw[->, line width=1pt] (0, 0) -- (0, 6) node[left] {Hourly Wage (\$)};

      % Grid
      \draw[gray!20] (1,0) grid (13, 5.5);

      % Male regression line
      \draw[thick, blue] (0, 1.5) -- (12, 4.8)
        node[above right, blue] {Male: $\beta_0 + \beta_1 \, \text{educ}$};

      % Female regression line (shifted down)
      \draw[thick, red] (0, 0.7) -- (12, 4.0)
        node[below right, red] {Female: $(\beta_0 + \delta_0) + \beta_1 \, \text{educ}$};

      % Illustrate intercept shift
      \draw[dashed, gray] (0, 1.5) -- (0, 0.7);
      \draw[<->, line width=1.5pt, green] (0.2, 1.5) -- (0.2, 0.7)
        node[left, green, font=\small] {$\delta_0 < 0$};

      % Annotate slopes
      \node[above, font=\small] at (3, 2.5) {Same slope: $\beta_1$};
    \end{tikzpicture}
  \end{center}

  \begin{keybox}
    If $\delta_0 < 0$, females earn \textit{less} than males with the same education.
    This difference is the \textbf{gender wage gap} (ceteris paribus).
  \end{keybox}

\end{frame}

\begin{frame}[shrink=25]{The Dummy Variable Trap}

  \textbf{Warning:} If you have $g$ categories, use \textbf{exactly} $g-1$ dummies!

  \medskip

  \textbf{Why?} Suppose we try to use:
  \[
    y = \beta_0 + \delta_{\text{male}} \cdot \textit{male} + \delta_{\text{female}} \cdot \textit{female} + u
  \]

  \textbf{Problem:} For every observation, either \texttt{male}=1, \texttt{female}=0 or \texttt{male}=0, \texttt{female}=1.

  This means: $\texttt{male} + \texttt{female} = 1$ (perfectly collinear with intercept $\beta_0$).

  \medskip

  \textbf{Matrix algebra:} The design matrix has linearly dependent columns:
  \[
    \begin{bmatrix} 1 & 0 & 1 \\ 1 & 1 & 0 \\ 1 & 1 & 0 \end{bmatrix}
    \quad \Rightarrow \quad
    \mathbf{X}^{\prime} \mathbf{X} \text{ is singular}
  \]

  Result: OLS cannot compute $(\mathbf{X}^{\prime} \mathbf{X})^{-1}$. Software will either:
  \begin{enumerate}
    \item Throw an error (singular matrix)
    \item Drop one dummy automatically (Stata, R do this)
  \end{enumerate}

  \medskip

  \begin{highlightbox}
    \textbf{Rule:} For $g$ categories, define $g-1$ dummies. The omitted category is the \textbf{base group} or \textbf{benchmark group}.

    All coefficients are interpreted \textit{relative to} the base group.
  \end{highlightbox}

  \medskip

  In our wage example: \texttt{female} is the only dummy; males (female=0) form the base group.

\end{frame}

\begin{frame}[shrink=30]{Example 7.1: Gender Wage Gap Regression}

  \textbf{Dataset:} WAGE1 (n = 526 workers)

  \vspace{-0.3cm}

  \begin{examplebox}{Wage Regression Results}
    \[
      \widehat{\textit{wage}} = -1.57 - 1.81 \, \textit{female} + 0.572 \, \textit{educ} + 0.25 \, \textit{exper} + 0.141 \, \textit{tenure}
    \]
    \[
      (0.72) \quad (0.26) \quad\quad\quad\; (0.049) \quad\quad (0.012) \quad\quad (0.021)
    \]
    \[
      n = 526, \quad R^2 = 0.364
    \]
  \end{examplebox}

  \medskip

  \textbf{Interpretation:}

  \begin{keybox}
    \begin{itemize}
      \item Holding education, experience, and tenure constant, women earn \textbf{\$1.81 per hour less} than men (in 1976 dollars).
      \item This is the \textbf{controlled gender wage gap} or \textbf{conditional} gap — the part not explained by these three factors.
      \item t-statistic for female: $t = -1.81 / 0.26 = -6.96$ (highly significant at any conventional level).
      \item \textbf{Economic magnitude:} At mean wage of \$5.58/hr, this represents a $1.81/5.58 = 32.4\%$ wage penalty.
    \end{itemize}
  \end{keybox}

  \medskip

  \begin{highlightbox}
    \textbf{Causality caveat:} Even with controls for educ, exper, tenure, causal interpretation requires $\mathbb{E}[u \mid \textit{female}, \textit{educ}, \textit{exper}, \textit{tenure}] = 0$.

    Omitted factors (ability, motivation, occupational choice, discrimination) bias the estimate.
  \end{highlightbox}

\end{frame}

\begin{frame}[shrink=25]{Controlled vs. Raw Gap}

  \textbf{Same regression, now estimate without controls:}

  \vspace{-0.3cm}

  \begin{center}
    \small
    \begin{tabular}{lcccc}
      \toprule
      & \textbf{Controlled} & & & \textbf{Raw} \\
      & (all controls) & & & (no controls) \\
      \midrule
      female & $-1.81$ & & & $-2.51$ \\
      & $(0.26)$ & & & $(0.30)$ \\
      \midrule
      & $n=526, R^2=0.364$ & & & $n=526, R^2=0.116$ \\
      \bottomrule
    \end{tabular}
  \end{center}

  \medskip

  \textbf{Why is the raw gap larger?}

  \begin{keybox}
    \begin{itemize}
      \item Raw gap (\$2.51) includes the effect of differences in education, experience, and tenure.
      \item Women in the sample have, on average, fewer years of work experience and lower tenure.
      \item Controlling for these factors \textit{reduces} the estimated wage gap to \$1.81.
      \item The difference $2.51 - 1.81 = 0.70$ is explained by \textit{compositional differences}.
    \end{itemize}
  \end{keybox}

\end{frame}

\begin{frame}[shrink=25]{Example 7.2: Computer Ownership and College GPA}

  \textbf{Question:} Does owning a personal computer improve college grades?

  \medskip

  \begin{examplebox}{PC Ownership and GPA (GPA1)}
    \[
      \widehat{\textit{colGPA}} = 1.26 + 0.157 \, \textit{PC} + 0.447 \, \textit{hsGPA} + 0.0087 \, \textit{ACT}
    \]
    \[
      (0.33) \quad (0.057) \quad\quad (0.094) \quad\quad\quad (0.0105)
    \]
    \[
      n = 141, \quad R^2 = 0.219
    \]
  \end{examplebox}

  \medskip

  \textbf{Interpretation:} PC ownership is associated with a $+0.157$ point higher GPA.

  \medskip

  \begin{keybox}
    \textbf{Causality concern:}
    \begin{itemize}
      \item Does owning a PC \textit{cause} higher GPA, or do stronger students \textit{buy} PCs?
      \item Possible omitted variable bias: ability, motivation, parental wealth.
      \item This is \textbf{observational data}, not a randomized experiment.
      \item Cannot conclude $\beta_{\text{PC}}$ is causal without additional assumptions.
    \end{itemize}
  \end{keybox}

\end{frame}

%% ============================================================
\section{\texorpdfstring{\S7-2a}{S7-2a} — Log(y) with Dummies}
%% ============================================================

\begin{frame}[shrink=20]{Interpreting Dummy Coefficients: Percentage Changes}

  When the \textbf{dependent variable is in logs}, a dummy coefficient has a special interpretation.

  \medskip

  \textbf{Approximate rule:} If $\hat{\beta}$ is small, $100 \hat{\beta} \approx \%$ change.

  \medskip

  \begin{definitionbox}{Exact Percentage Interpretation}
    \[
      \log(y) = \beta_0 + \delta \cdot d + \text{other terms}
    \]

    When $d$ changes from 0 to 1:
    \[
      \% \text{ change in } y = 100 \times \left[ \exp(\hat{\delta}) - 1 \right] \quad \text{(eq.\ 7.10)}
    \]

    If $\hat{\delta}$ is small (e.g., $< 0.10$), then $100 \hat{\delta} \approx \%$ change.
  \end{definitionbox}

  \medskip

  \textbf{Example:} If $\hat{\delta} = 0.054$:
  \begin{itemize}
    \item Approximate: $100 \times 0.054 = 5.4\%$ \\
    \item Exact: $100 \times (\exp(0.054) - 1) = 100 \times 0.0555 = 5.55\%$
  \end{itemize}

\end{frame}

\begin{frame}[shrink=25]{Example 7.4: Housing Prices and Colonial Style}

  \begin{examplebox}{Housing Price Regression (HPRICE1)}
    \[
      \widehat{\log(\textit{price})} = -1.35 + 0.168 \log(\textit{lotsize}) + 0.707 \log(\textit{sqrft})
    \]
    \[
      (0.65) \quad (0.038) \quad\quad\quad\quad\quad (0.093)
    \]
    \[
      + 0.027 \, \textit{bdrms} + 0.054 \, \textit{colonial}
    \]
    \[
      \quad\quad (0.029) \quad\quad\quad (0.045)
    \]
    \[
      n = 88, \quad R^2 = 0.649
    \]
  \end{examplebox}

  \medskip

  \textbf{Interpretation of colonial dummy:}

  \begin{keybox}
    \begin{itemize}
      \item Approximate: Colonial homes sell for $5.4\%$ more.
      \item Exact: $\exp(0.054) - 1 = 0.0555 \Rightarrow 5.55\%$ more.
      \item t-statistic: $0.054 / 0.045 = 1.2$ (not significant at 5\%).
    \end{itemize}
  \end{keybox}

\end{frame}

\begin{frame}[shrink=25]{Example 7.5: Log Hourly Wage and Gender}

  \textbf{Same dataset (WAGE1) as Example 7.1, but with log(wage) as dependent variable.}

  \medskip

  \begin{examplebox}{Log Wage with Gender Dummy}
    \[
      \widehat{\log(\textit{wage})} = 0.417 - 0.297 \, \textit{female} + 0.080 \, \textit{educ} + 0.029 \, \textit{exper}
    \]
    \[
      (0.099) \quad (0.036) \quad\quad (0.007) \quad\quad (0.005)
    \]
    \[
      - 0.00058 \, \textit{exper}^2 + 0.032 \, \textit{tenure} - 0.00059 \, \textit{tenure}^2
    \]
    \[
      n = 526, \quad R^2 = 0.441
    \]
  \end{examplebox}

  \medskip

  \textbf{Interpretation of female coefficient:}

  \begin{keybox}
    \begin{itemize}
      \item Approximate: Women earn $-29.7\%$ less than men.
      \item Exact: $\exp(-0.297) - 1 = -0.257 \Rightarrow -25.7\%$ less.
      \item The log model sometimes gives large coeff. estimates relative to % interpretation.
    \end{itemize}
  \end{keybox}

\end{frame}

%% ============================================================
\section{Using Dummy Variables for Multiple Categories (\S7-3)}
%% ============================================================

\begin{frame}[shrink=20]{Multiple Categories: The g-1 Rule}

  \textbf{General problem:} Encode a categorical variable with $g$ categories.

  \medskip

  \textbf{Solution:} Define \textbf{exactly} $g-1$ dummies. The omitted category is the reference.

  \medskip

  \begin{highlightbox}
    \textbf{Example:} Education with 4 levels: high school (base), some college, college degree, graduate degree.

    Define 3 dummies:
    \begin{itemize}
      \item $d_1 = 1$ if some college, 0 otherwise
      \item $d_2 = 1$ if college degree, 0 otherwise
      \item $d_3 = 1$ if graduate degree, 0 otherwise
    \end{itemize}

    All coefficients are relative to high school (the base group).
  \end{highlightbox}

  \medskip

  \begin{keybox}
    \textbf{To identify the base group:} Look at the regression output. The category with \textit{no} dummy variable is the reference.
  \end{keybox}

\end{frame}

\begin{frame}[shrink=28]{Example 7.6: Four Gender-Marital Status Groups}

  \textbf{Base group:} Single males. Define three dummies:
  \begin{itemize}
    \item \texttt{marrmale} = 1 if married male
    \item \texttt{marrfem} = 1 if married female
    \item \texttt{singfem} = 1 if single female
  \end{itemize}

  \medskip

  \begin{examplebox}{Four Gender-Marital Groups (WAGE1)}
    \[
      \widehat{\log(\textit{wage})} = 0.321 + 0.213 \, \textit{marrmale} - 0.198 \, \textit{marrfem}
    \]
    \[
      (0.100) \quad (0.055) \quad\quad\quad\quad (0.058)
    \]
    \[
      - 0.110 \, \textit{singfem} + 0.079 \, \textit{educ} + 0.027 \, \textit{exper} - 0.00054 \, \textit{exper}^2
    \]
    \[
      + 0.029 \, \textit{tenure} - 0.00053 \, \textit{tenure}^2 \quad\quad n = 526, R^2 = 0.461
    \]
  \end{examplebox}

  \medskip

  \begin{keybox}
    \begin{itemize}
      \item Married men earn $+21.3\%$ more than single men (base).
      \item Married women earn $-19.8\%$ less than single men.
      \item Single women earn $-11.0\%$ less than single men.
    \end{itemize}
  \end{keybox}

\end{frame}

\begin{frame}[shrink=25]{Ordinal Variables and Credit Ratings}

  \textbf{Question:} How does credit rating (ordinal) affect mortgage approval probability?

  \medskip

  \textbf{Credit ratings:} 0 (no credit), 1 (poor), 2 (fair), 3 (good), 4 (excellent).

  \medskip

  \begin{definitionbox}{Ordinal Dummy Model}
    \[
      \textit{approve} = \beta_0 + \delta_1 \, \textit{CR1} + \delta_2 \, \textit{CR2} + \delta_3 \, \textit{CR3} + \delta_4 \, \textit{CR4} + u
    \]

    Base group: CR=0 (no credit history).
  \end{definitionbox}

  \medskip

  \textbf{Key insight:} By treating each category separately (rather than using CR as a continuous score), we allow the effect to be nonlinear.

  \medskip

  \begin{keybox}
    The coefficients $\delta_1, \delta_2, \delta_3, \delta_4$ reveal the exact shape of the relationship.
    If they increase monotonically, the effect is monotonic.
  \end{keybox}

\end{frame}

\begin{frame}[shrink=28]{Example 7.8: Law School Rankings and Salaries}

  \textbf{Question:} How much does a law school's ranking matter for starting salary?

  \medskip

  \begin{examplebox}{Law School Rankings and Salaries}
    \begin{small}
    \[
      \widehat{\log(\textit{salary})} = 9.17 + 0.700 \, \textit{top10} + 0.594 \, \textit{r11\_25}
    \]
    \[
      + 0.375 \, \textit{r26\_40} + 0.263 \, \textit{r41\_60} + 0.132 \, \textit{r61\_100}
    \]
    \[
      + 0.0057 \, \textit{LSAT} + 0.041 \, \textit{GPA} + \text{(other controls)}
    \]
    \[
      n = 136, \quad R^2 = 0.911, \quad \bar{R}^2 = 0.905
    \]
    \end{small}
  \end{examplebox}

  \medskip

  \textbf{Interpretation:}

  \begin{keybox}
    \begin{itemize}
      \item Base group: schools ranked outside top 100 (median salary).
      \item Top-10 schools: $\exp(0.700) - 1 \approx +101\%$ salary premium (median salary roughly doubles!).
      \item Ranks 11-25: $\exp(0.594) - 1 \approx +81\%$ premium.
      \item Ranks 26-40: $\exp(0.375) - 1 \approx +46\%$ premium.
      \item Clear rank order: effect diminishes as ranking worsens (monotonic).
    \end{itemize}
  \end{keybox}

  \medskip

  \begin{highlightbox}
    \textbf{Economic significance:} School ranking is the dominant determinant of starting salary (other variables like LSAT and GPA are much smaller). Controls for ability (LSAT, GPA) don't eliminate ranking effects → suggests \textit{signaling} or \textit{network effects}, not just ability.
  \end{highlightbox}

\end{frame}

%% ============================================================
\section{Interactions Involving Dummy Variables (\S7-4)}
%% ============================================================

\begin{frame}[shrink=25]{Interactions Among Dummy Variables}

  \textbf{Motivation:} Do different groups respond differently to the same factor?

  \medskip

  \textbf{Example:} Does marriage have the same wage effect for men and women?

  \medskip

  \begin{definitionbox}{Dummy × Dummy Interaction}
    \[
      \log(\textit{wage}) = \beta_0 + \delta_0 \textit{female} + \delta_1 \textit{married}
    \]
    \[
      + \gamma \, (\textit{female} \times \textit{married}) + \text{other controls} + u
    \]
  \end{definitionbox}

  \medskip

  \textbf{Interpretation:} The coefficient $\gamma$ captures the \textit{additional} effect of marriage for women.

  \begin{keybox}
    \begin{itemize}
      \item For men: marriage raises wage by $\delta_1$ (alone).
      \item For women: marriage raises wage by $\delta_1 + \gamma$.
      \item If $\gamma < 0$, marriage "helps" men more than women (differential return).
    \end{itemize}
  \end{keybox}

\end{frame}

\begin{frame}[shrink=25]{Example 7.9: Computer Use at Work vs. Home}

  \textbf{Dataset:} 1989 CPS microdata (n = 13,379 workers); Krueger (1993)

  \medskip

  \begin{examplebox}{Computer Use and Wages}
    \[
      \widehat{\log(\textit{wage})} = \hat{\beta}_0 + 0.177 \, \textit{compwork} + 0.070 \, \textit{comphome}
    \]
    \[
      \quad\quad\quad\quad\quad\quad (0.009) \quad\quad\quad\quad (0.019)
    \]
    \[
      + 0.017 \, (\textit{compwork} \times \textit{comphome}) + \text{other controls}
    \]
    \[
      \quad\quad\quad (0.023)
    \]
    \smallskip

    Data: 1989 CPS microdata (n = 13,379)
  \end{examplebox}

  \medskip

  \textbf{Interpretation:}

  \begin{keybox}
    \begin{itemize}
      \item Using a computer at work: $+17.7\%$ wage premium (large and highly significant).
      \item Using at home (no work): $+7.0\%$ premium.
      \item Interaction (using both): $+1.7\%$ additional (not significant, t $\approx 0.74$).
      \item Two main findings: (1) workplace computer use valuable, (2) home use has modest effect.
    \end{itemize}
  \end{keybox}

\end{frame}

\begin{frame}[shrink=20]{Allowing Different Slopes: Gender × Education}

  \textbf{Question:} Do men and women get different returns to education?

  \medskip

  \begin{definitionbox}{Dummy × Continuous Interaction}
    \[
      \log(\textit{wage}) = \beta_0 + \delta_0 \textit{female} + \beta_1 \textit{educ}
    \]
    \[
      + \delta_1 \, (\textit{female} \times \textit{educ}) + \text{other controls} + u
    \]
  \end{definitionbox}

  \medskip

  \textbf{Interpretation:}
  \begin{itemize}
    \item $\beta_1$ = return to education for \textbf{males}.
    \item $\delta_1$ = \textbf{additional} return to education for females (compared to males).
    \item Total return for females: $\beta_1 + \delta_1$.
  \end{itemize}

  \medskip

  \begin{keybox}
    If $\delta_1 \neq 0$ (and significant), the market \textit{values} education differently for men vs. women.
    This is a type of \textbf{discrimination by skill} or \textbf{skill-biased discrimination}.
  \end{keybox}

\end{frame}

\begin{frame}[shrink=15]{Four Cases of Slope Shifts}

  \begin{center}
    \begin{tikzpicture}[scale=0.9]
      % Top-left: delta0 > 0, delta1 < 0 (intercept up, slope down)
      \begin{scope}[xshift=-7cm, yshift=2.5cm]
        \draw[->, line width=1pt] (0, 0) -- (4, 0) node[below] {$x$};
        \draw[->, line width=1pt] (0, 0) -- (0, 3.5) node[left] {$y$};
        \draw[thick, blue] (0, 1.5) -- (4, 3.2) node[above right, blue, font=\small] {Males};
        \draw[thick, red] (0, 2.0) -- (4, 1.8) node[below right, red, font=\small] {Females};
        \node[below] at (2, -0.8) {$\delta_0 > 0, \delta_1 < 0$};
      \end{scope}

      % Top-right: delta0 > 0, delta1 > 0 (both up)
      \begin{scope}[xshift=0cm, yshift=2.5cm]
        \draw[->, line width=1pt] (0, 0) -- (4, 0) node[below] {$x$};
        \draw[->, line width=1pt] (0, 0) -- (0, 3.5) node[left] {$y$};
        \draw[thick, blue] (0, 1.5) -- (4, 3.2) node[above right, blue, font=\small] {Males};
        \draw[thick, red] (0, 2.0) -- (4, 4.0) node[above right, red, font=\small] {Females};
        \node[below] at (2, -0.8) {$\delta_0 > 0, \delta_1 > 0$};
      \end{scope}

      % Bottom-left: delta0 < 0, delta1 < 0 (both down)
      \begin{scope}[xshift=-7cm, yshift=-2.5cm]
        \draw[->, line width=1pt] (0, 0) -- (4, 0) node[below] {$x$};
        \draw[->, line width=1pt] (0, 0) -- (0, 3.5) node[left] {$y$};
        \draw[thick, blue] (0, 2.5) -- (4, 4.2) node[above right, blue, font=\small] {Males};
        \draw[thick, red] (0, 1.0) -- (4, 1.8) node[below right, red, font=\small] {Females};
        \node[below] at (2, -0.8) {$\delta_0 < 0, \delta_1 < 0$};
      \end{scope}

      % Bottom-right: delta0 < 0, delta1 > 0 (intercept down, slope up)
      \begin{scope}[xshift=0cm, yshift=-2.5cm]
        \draw[->, line width=1pt] (0, 0) -- (4, 0) node[below] {$x$};
        \draw[->, line width=1pt] (0, 0) -- (0, 3.5) node[left] {$y$};
        \draw[thick, blue] (0, 2.5) -- (4, 3.2) node[above right, blue, font=\small] {Males};
        \draw[thick, red] (0, 1.0) -- (4, 4.0) node[above right, red, font=\small] {Females};
        \node[below] at (2, -0.8) {$\delta_0 < 0, \delta_1 > 0$};
      \end{scope}
    \end{tikzpicture}
  \end{center}

  \textbf{In each case:} females (red line) have a different intercept (\textit{offset}) and/or a different slope compared to males (blue).

\end{frame}

\begin{frame}[shrink=28]{Example 7.10: Gender and Return to Education}

  \begin{examplebox}{Gender-Education Interaction}
    Log wage regression with WAGE1 data:
    \[
      \widehat{\log(\textit{wage})} = 0.389 - 0.227 \, \textit{female} + 0.082 \, \textit{educ}
    \]
    \[
      (0.119) \quad (0.168) \quad\quad\quad (0.008)
    \]
    \[
      - 0.0056 \, (\textit{female} \times \textit{educ}) + 0.029 \, \textit{exper} - 0.00058 \, \textit{exper}^2
    \]
    \[
      (0.0131) \quad\quad\quad\quad\quad\quad\quad (0.005) \quad\quad (0.00011)
    \]
    \[
      + 0.032 \, \textit{tenure} - 0.00059 \, \textit{tenure}^2
    \]
    \[
      n = 526, \quad R^2 = 0.411
    \]
  \end{examplebox}

  \medskip

  \textbf{Interpretation:}

  \begin{keybox}
    \begin{itemize}
      \item Return to education for \textbf{males}: $0.82\%$ per year ($t = 10.25$, highly significant).
      \item Return for \textbf{females}: $0.82 - 0.56 = 0.76\%$ per year.
      \item Difference is \textbf{not statistically significant} ($t = -0.0056 / 0.0131 = -0.43$).
      \item \textbf{Warning:} Large SE on female coeff. (0.168) due to high correlation between \texttt{female} and \texttt{female}×\texttt{educ}.
    \end{itemize}
  \end{keybox}

\end{frame}

\begin{frame}[shrink=25]{The Chow Test: Comparing Regression Functions}

  \textbf{Question:} Do two groups have fundamentally \textit{different} regression relationships?

  \medskip

  \textbf{Setup:} Compare a \textbf{pooled} model vs. separate models by group.

  \medskip

  \begin{highlightbox}
    \textbf{Chow Test Statistic (eq.\ 7.24):}
    \[
      F = \frac{(\text{SSR}_P - (\text{SSR}_1 + \text{SSR}_2))/(k+1)}{(\text{SSR}_1 + \text{SSR}_2)/(n - 2(k+1))}
    \]

    where:
    \begin{itemize}
      \item $\text{SSR}_P$ = sum of squared residuals from pooled (unrestricted) model
      \item $\text{SSR}_1, \text{SSR}_2$ = SSR from separate regressions for each group
      \item $k$ = number of slopes (plus intercept $\Rightarrow$ $k+1$ parameters per group)
      \item Numerator: improvement in fit from allowing different intercepts and slopes
    \end{itemize}
  \end{highlightbox}

  \medskip

  \textbf{Null hypothesis:} $H_0: \delta_0 = \delta_1 = \cdots = \delta_k = 0$ (both groups follow the same regression).

  \textbf{Under} $H_0$: $F \sim F_{k+1, n-2(k+1)}$ (F distribution with $k+1$ and $n-2(k+1)$ degrees of freedom).

  \medskip

  \begin{keybox}
    \textbf{Intuition:} The Chow test asks: ``Is the pooled model (restriction) significantly worse than the group-specific models (unrestricted)?''

    If $F$ is large, the groups are \textbf{very different} (different intercepts and/or slopes).

    If $F$ is small, a single regression fits both groups well.
  \end{keybox}

  \medskip

  \textbf{Two variants:}
  \begin{itemize}
    \item \textbf{Full Chow test:} Allow both intercepts and all slopes to differ ($k+1$ restrictions)
    \item \textbf{Restricted Chow:} Allow only intercepts to differ; slopes same ($1$ restriction: test $\delta_0 = 0$ only)
  \end{itemize}

\end{frame}

%% ============================================================
\section{A Binary Dependent Variable: The Linear Probability Model (\S7-5)}
%% ============================================================

\begin{frame}[shrink=20]{Binary Dependent Variables: Motivation}

  \textbf{New challenge:} What if the dependent variable is binary (0 or 1)?

  \medskip

  \begin{examplebox}{Examples of Binary Outcomes}
    \begin{itemize}
      \item Criminology: arrested or not arrested (1 = arrested)
      \item Labor economics: employed or not (1 = employed)
      \item Urban economics: home ownership: yes or no (1 = owner)
      \item Education: college attendance: yes or no (1 = attended)
      \item Finance: loan approval: approved or denied (1 = approved)
      \item Health: disease present or absent (1 = present)
    \end{itemize}
  \end{examplebox}

  \medskip

  \textbf{Question:} Can we still use OLS regression?

  \medskip

  \begin{highlightbox}
    Yes! OLS still gives us \textbf{unbiased} estimates of the mean effect (MLR.1–4 sufficient).

    The predicted value $\hat{y}$ can be interpreted as the \textbf{probability of success} given $x$.
  \end{highlightbox}

  \medskip

  \begin{keybox}
    But OLS has \textbf{three key limitations} when $y$ is binary. Understanding these motivates nonlinear methods (Probit, Logit — Chapter 17).
  \end{keybox}

\end{frame}

\begin{frame}[shrink=25]{The Linear Probability Model (LPM)}

  \begin{definitionbox}{Linear Probability Model (LPM)}
    \[
      y = \beta_0 + \beta_1 x_1 + \cdots + \beta_k x_k + u \quad \text{(eq.\ 7.27)}
    \]

    where $y \in \{0, 1\}$ (binary outcome).

    OLS estimates the \textbf{conditional probability}:
    \[
      \E[y | x] = P(y = 1 | x) = \beta_0 + \beta_1 x_1 + \cdots + \beta_k x_k
    \]

    This assumes $u \in \{-p(x), 1-p(x)\}$ where $p(x) = \beta_0 + \beta_1 x_1 + \cdots + \beta_k x_k$.
  \end{definitionbox}

  \medskip

  \textbf{Slope interpretation:} $\beta_j$ is the marginal effect on the probability of success.

  \medskip

  \begin{highlightbox}
    \[
      \Delta P(y=1|x) = \beta_j \, \Delta x_j \quad \text{(eq.\ 7.28)}
    \]

    \textbf{Concrete example:} If $\beta_j = 0.05$ and $x_j$ increases by 1 unit, the probability of success \textit{increases by 5 percentage points}.

    Compare to: if $y$ were continuous wage and $\beta_j = 0.05$, wage increases by \$0.05 per unit.
  \end{highlightbox}

\end{frame}

\begin{frame}[shrink=20]{LPM Drawbacks}

  Despite its simplicity, the LPM has three key limitations:

  \medskip

  \begin{enumerate}
    \item \textbf{Predicted probabilities can fall outside [0, 1] (nonsensical predictions).}
    \begin{itemize}
      \item Example: If $\hat{y} = 0.2 - 0.15 \times \text{age}$ and age = 5, then $\hat{y} = 0.2 - 0.75 = -0.55$ (impossible probability!).
      \item Or if $\hat{\beta}_0 = 1.5$ and $\hat{\beta}_1 = 0.1$, at age = 100, $\hat{y} = 1.5 + 10 = 11.5$ (probability $> 1$).
      \item Nonlinear models (Probit, Logit) constrain predictions to [0, 1] via sigmoid or CDF function.
    \end{itemize}

    \vspace{0.3cm}

    \item \textbf{Constant marginal effects are unrealistic.}
    \begin{itemize}
      \item LPM assumes $\partial P(y=1) / \partial x = \beta_j$ everywhere (constant line).
      \item Reality: effects should flatten at extremes (diminishing returns).
      \item Hard to go from 95\% success to 100\% by increasing $x$ (probability ceiling) vs. 50\% to 65\% (linear zone).
    \end{itemize}

    \vspace{0.3cm}

    \item \textbf{Inherent heteroskedasticity violates MLR.5:}
    \[
      \Var(u | x) = p(x) \cdot [1 - p(x)] \quad \text{(eq.\ 7.30)}
    \]
    \begin{itemize}
      \item Variance is \textbf{largest} when $p(x) = 0.5$ (most uncertainty); \textbf{smallest} at extremes ($p \approx 0$ or 1).
      \item OLS standard errors computed assuming $\Var(u) = \sigma^2$ (constant) are \textbf{incorrect}.
      \item \textbf{Solution:} Use heteroskedasticity-robust standard errors (covered in Lecture 7).
    \end{itemize}
  \end{enumerate}

\end{frame}

\begin{frame}[shrink=30]{Example 7.12: Linear Probability Model of Arrests}

  \textbf{Dataset:} CRIME1 (n = 2,725 men)

  \textbf{Dependent variable:} \texttt{arr86} = 1 if arrested in 1986, 0 otherwise

  \medskip

  \begin{examplebox}{Arrest Probability Model}
    \[
      \widehat{\textit{arr86}} = 0.441 - 0.162 \, \textit{pcnv} + 0.0061 \, \textit{avgsen}
    \]
    \[
      (0.017) \quad (0.021) \quad\quad\quad (0.0065)
    \]
    \[
      - 0.0023 \, \textit{tottime} - 0.022 \, \textit{ptime86} - 0.043 \, \textit{qemp86}
    \]
    \[
      (0.0050) \quad\quad (0.005) \quad\quad (0.005)
    \]
    \[
      R^2 = 0.0474, \quad n = 2725
    \]
  \end{examplebox}

  \medskip

  \textbf{Interpretation:}

  \begin{keybox}
    \begin{itemize}
      \item Prior conviction (pcnv): reduces arrest probability by $-16.2$ percentage points (large effect).
      \item Quarter employed (qemp86): each additional quarter employed reduces arrest probability by $-4.3$ percentage points.
      \item Predicted $\hat{y}$ ranges from 0.04 to 0.86 (mostly in [0, 1], but some edge cases).
    \end{itemize}
  \end{keybox}

\end{frame}

\begin{frame}[shrink=20]{When LPM Works Well}

  Despite its drawbacks, LPM is useful when:

  \medskip

  \begin{highlightbox}
    \begin{itemize}
      \item You have a large sample (heteroskedasticity less critical with robust SEs).
      \item Predicted probabilities cluster away from 0 and 1 boundaries.
      \item Your goal is \textbf{estimation and inference}, not prediction.
      \item You need a simple, interpretable model with no distributional assumptions.
    \end{itemize}
  \end{highlightbox}

  \medskip

  \begin{keybox}
    \textbf{Standard practice:} Report LPM results, but always use robust standard errors to account for heteroskedasticity.

    For comparison: Lecture 15 (Limited Dependent Variables) covers nonlinear alternatives (Probit, Logit, Tobit).
  \end{keybox}

\end{frame}

%% ============================================================
\section{Policy Analysis and Program Evaluation (\S7-6)}
%% ============================================================

\begin{frame}[shrink=20]{Program Evaluation: The Fundamental Problem}

  \textbf{Motivating question:} Does a job training program increase earnings?

  \medskip

  \begin{examplebox}{Setting}
    \begin{itemize}
      \item Treatment group: workers who receive training.
      \item Control group: workers who don't.
      \item Goal: estimate the \textbf{causal effect} of training on future earnings.
    \end{itemize}
  \end{examplebox}

  \medskip

  \textbf{Challenge:} \textbf{Self-selection bias} (fundamental identification problem).

  \begin{keybox}
    \begin{itemize}
      \item Workers who volunteer for training may be more \textbf{motivated}, have higher \textbf{ability}, or expect greater \textbf{returns}.
      \item Simply comparing trained vs. untrained confounds:
      \[
        \mathbb{E}[\text{earn} \mid \text{trained}] - \mathbb{E}[\text{earn} \mid \text{untrained}]
      \]
      \[
        = \underbrace{\text{Training effect}}_{\text{what we want}} + \underbrace{\text{Selection bias}}_{\text{confounding}}
      \]
      \item Regression adjustment (conditioning on observables) can \textit{partially} control for selection.
      \item But if training choice depends on unobservables (motivation, ability, expectations), bias remains.
    \end{itemize}
  \end{keybox}

\end{frame}

\begin{frame}[shrink=20]{Restricted Regression Adjustment (RRA)}

  \begin{definitionbox}{RRA Model}
    \[
      y_i = \beta_0 + \tau \, w_i + \beta_1 x_{i1} + \cdots + \beta_k x_{ik} + u_i
    \]

    where $w_i = 1$ if worker $i$ received training, 0 otherwise;
    $x_1, \ldots, x_k$ are control variables (age, education, prior earnings, etc.);
    $\tau$ is the treatment effect.
  \end{definitionbox}

  \medskip

  \textbf{Interpretation:} The coefficient $\hat{\tau}$ is the difference in earnings between trained and untrained workers, \textit{holding controls constant}.

  \medskip

  \begin{keybox}
    \textbf{Assumption:} All variables affecting earnings that differ between trained and untrained are \textit{observed} and included as controls.

    Any unobserved differences bias the estimate (omitted variable bias).
  \end{keybox}

\end{frame}

\begin{frame}[shrink=20]{Unrestricted Regression Adjustment (URA)}

  \textbf{Question:} Does training effect differ by worker characteristics (age, education)?

  \medskip

  \begin{definitionbox}{URA with Demeaned Interactions}
    \[
      y_i = \beta_0 + \tau \, w_i + \sum_j \beta_j (x_{ij} - \bar{x}_j)
    \]
    \[
      + \sum_j \gamma_j \, w_i \cdot (x_{ij} - \bar{x}_j) + u_i
    \]

    where $(x_{ij} - \bar{x}_j)$ are \textbf{demeaned} controls.
  \end{definitionbox}

  \medskip

  \textbf{Key insight:} Demeaning allows $\tau$ to be the \textbf{average treatment effect (ATE)} at the sample means of $x_1, \ldots, x_k$.

  \medskip

  \begin{highlightbox}
    The interaction coefficients $\gamma_j$ reveal how the training effect varies by characteristics.

    Example: If $\gamma_{\text{age}} > 0$, older workers benefit more from training.
  \end{highlightbox}

\end{frame}

\begin{frame}[shrink=30]{Example 7.13: Job Training Program Evaluation}

  \textbf{Dataset:} JTRAIN98 (n = 1,130 workers); outcome = earnings in 1998

  \medskip

  \begin{small}
  \begin{examplebox}{Job Training Evaluation}
    \[
      \widehat{\textit{earn98}} = 5.08 + 3.11 \, \textit{train} + 0.353 \, (\textit{earn96} - \bar{e}_{96})
    \]
    \[
      + 0.378 \, (\textit{educ} - \bar{e}) - 0.196 \, (\textit{age} - \bar{a})
    \]
    \[
      + 0.133 \, \textit{train} \cdot (\textit{earn96} - \bar{e}_{96}) - 0.035 \, \textit{train} \cdot (\textit{educ} - \bar{e})
    \]
    \[
      + 0.508 \, \textit{train} \cdot (\textit{age} - \bar{a}) - 0.993 \, \textit{train} \cdot (\textit{married} - \bar{m})
    \]
    \[
      n = 1130, \quad R^2 = 0.409
    \]
  \end{examplebox}
  \end{small}

  \medskip

  \begin{keybox}
    \begin{itemize}
      \item \textbf{Average treatment effect:} $\hat{\tau} = 3.11$ thousand dollars (se = 0.53).
      \item Training increases expected earnings by about \$3,110.
      \item \textbf{Heterogeneous effects:} Training benefit \textit{increases} with prior earnings and age, but \textit{decreases} with education and marriage.
    \end{itemize}
  \end{keybox}

\end{frame}

%% ============================================================
\section{Interpreting Regression with Discrete Dependent Variables (\S7-7)}
%% ============================================================

\begin{frame}[shrink=20]{Discrete Dependent Variables: Broad Interpretation}

  \textbf{Key principle:} OLS estimates $\E[y | x]$, whether $y$ is continuous or discrete.

  \medskip

  \begin{highlightbox}
    When $y$ is discrete (e.g., number of children, years of education, job rank), OLS coefficients still have a meaningful interpretation:

    \[
      \beta_j = \frac{\partial \E[y | x]}{\partial x_j}
    \]

    This is the change in \textbf{expected value} of $y$ when $x_j$ increases.
  \end{highlightbox}

  \medskip

  \textbf{Example interpretation:} "Each additional year of education reduces the expected number of children by 0.09."

  \medskip

  \begin{keybox}
    You do \textit{not} need Probit, Logit, or other nonlinear methods for discrete $y$.
    OLS still works, as long as you interpret coefficients as effects on the \textit{mean} of $y$.
  \end{keybox}

\end{frame}

\begin{frame}[shrink=25]{Example: Fertility and Education}

  \textbf{Dataset:} FERTIL2 (n = 4,361 women)

  \medskip

  \begin{examplebox}{Fertility and Education}
    \[
      \widehat{\textit{children}} = -1.997 + 0.175 \, \textit{age} - 0.090 \, \textit{educ} + \ldots
    \]
    \[
      n = 4361, \quad R^2 = 0.560 \quad \text{(eq.\ 7.45)}
    \]
  \end{examplebox}

  \medskip

  \textbf{Interpretation of education coefficient:}

  \begin{keybox}
    Each additional year of education is associated with $-0.090$ fewer children on average.

    This is a \textit{large} effect: 12 more years of education → $12 \times 0.090 = 1.08$ fewer children.
  \end{keybox}

  \medskip

  \textbf{Adding a dummy variable:}

  \medskip

  \begin{examplebox}{With Electric Dummy}
    \[
      \widehat{\textit{children}} = \ldots - 0.362 \, \textit{electric} + \ldots
    \]

    Having electricity in the home is associated with 0.36 fewer children (eq.\ 7.47).
  \end{examplebox}

\end{frame}

%% ============================================================
\section{Summary and Key Takeaways}
%% ============================================================

\begin{frame}[shrink=25]{Key Concepts Summary}

  \begin{keybox}
    \textbf{Dummy Variables:}
    \begin{itemize}
      \item Encode qualitative info as 0/1; name clearly (female, not gender).
      \item For $g$ categories, use $g-1$ dummies; omitted category = base group.
      \item All coefficients interpreted relative to base group.
    \end{itemize}
  \end{keybox}

  \medskip

  \begin{highlightbox}
    \textbf{Log(y) with Dummies:}
    \begin{itemize}
      \item Approximate: $100 \hat{\beta} \approx \%$ change (if $\hat{\beta}$ is small).
      \item Exact: $100 \times [\exp(\hat{\beta}) - 1] = \%$ change.
    \end{itemize}
  \end{highlightbox}

  \medskip

  \begin{keybox}
    \textbf{Interactions:}
    \begin{itemize}
      \item Dummy × dummy: allows different intercepts by group.
      \item Dummy × continuous: allows different slopes by group.
      \item Chow test: formally test if two groups have different regression functions.
    \end{itemize}
  \end{keybox}

\end{frame}

\begin{frame}[shrink=20]{Key Concepts (continued)}

  \begin{highlightbox}
    \textbf{Linear Probability Model (Binary y):}
    \begin{itemize}
      \item $\hat{y}$ is predicted probability; $\beta_j$ = change in probability per unit of $x_j$.
      \item Drawbacks: predictions can exceed [0, 1]; constant marginal effects; inherent heteroskedasticity.
      \item Solution: always use robust standard errors (account for Var$(y|x) = p(1-p)$).
    \end{itemize}
  \end{highlightbox}

  \medskip

  \begin{keybox}
    \textbf{Program Evaluation:}
    \begin{itemize}
      \item RRA: simple regression adjustment (assumes no unobserved confounders).
      \item URA: allows treatment effect to vary by controls (demeaned interactions).
      \item Both require assumption that all confounders are observed.
    \end{itemize}
  \end{keybox}

  \medskip

  \begin{examplebox}{Discrete Dependent Variables}
    OLS is valid even when $y$ is discrete (0/1, count, rank). Coefficients measure effects on $\E[y | x]$.
  \end{examplebox}

\end{frame}

\begin{frame}[shrink=25]{Bridge to Lecture 7: Heteroskedasticity}

  \textbf{Why Chapter 8 matters:}

  The Linear Probability Model revealed an important limitation of OLS:

  \vspace{-0.3cm}

  \begin{highlightbox}
    When $y$ is binary: $\Var(u | x) = p(x) \cdot [1 - p(x)]$ — the error variance \textit{depends on} $x$.

    This violates MLR.5 (\textbf{Homoskedasticity}).
  \end{highlightbox}

  \medskip

  \textbf{Broader problem:} Many real datasets exhibit \textbf{heteroskedasticity}:

  \begin{examplebox}{Examples}
    \begin{itemize}
      \item Household consumption: variance increases with income.
      \item Firm profit: variance increases with firm size.
      \item Test scores: variance decreases as study time increases.
    \end{itemize}
  \end{examplebox}

  \medskip

  \begin{keybox}
    When heteroskedasticity is present:
    \begin{itemize}
      \item OLS estimates are still \textbf{unbiased} and \textbf{consistent}.
      \item But standard errors are \textbf{wrong} (too small or too large).
      \item Solution: \textbf{Robust standard errors} (Lecture 7).
    \end{itemize}
  \end{keybox}

\end{frame}

%% ============================================================
\end{document}
